% !TeX root = ../sections/03_solutions.tex
\subsection{1.1.B.5}
\begin{proof}
    求める数列を
    \[
        x_n=
        \frac{a_1+2a_2+\cdots+(n-1)a_{n-1}+na_n}{n^2}
    \]
    とおく。

    シグマ記号を用いると，
    \[
        x_n=
        \frac{1}{n^2}\sum_{k=1}^{n}ka_k
    \]
    である。

    仮定より
    \[
        \lim_{k\to\infty}a_k=\alpha
    \]
    であるから，
    \[
        a_k=\alpha+(a_k-\alpha)
    \]
    と分解する。

    これを代入すると，
    \[
        x_n
        =
        \frac{1}{n^2}\sum_{k=1}^{n}k\{\alpha+(a_k-\alpha)\}
    \]
    である。したがって，
    \[
        x_n
        =
        \alpha\frac{1}{n^2}\sum_{k=1}^{n}k
        +
        \frac{1}{n^2}\sum_{k=1}^{n}k(a_k-\alpha)
    \]
    となる。

    ここで
    \[
        \sum_{k=1}^{n}k=\frac{n(n+1)}{2}
    \]
    であるから，
    \[
        \alpha\frac{1}{n^2}\sum_{k=1}^{n}k
        =
        \alpha\frac{n(n+1)}{2n^2}
        =
        \frac{\alpha}{2}\left(1+\frac{1}{n}\right)
    \]
    である。

    よって
    \[
        \alpha\frac{1}{n^2}\sum_{k=1}^{n}k
        \to
        \frac{\alpha}{2}
    \]
    である。

    したがって，あとは
    \[
        \frac{1}{n^2}\sum_{k=1}^{n}k(a_k-\alpha)
        \to 0
    \]
    を示せばよい。
    そこで
    \[
        R_n=
        \frac{1}{n^2}\sum_{k=1}^{n}k(a_k-\alpha)
    \]
    とおく。

    $R_n\to 0$ を示すため，任意に $\epsilon>0$ を取る。

    $a_k\to\alpha$ であるから，
    \[
        \frac{\epsilon}{2}>0
    \]
    に対して，ある自然数 $N_1$ が存在して，
    \[
        k\geq N_1
    \]
    ならば
    \[
        |a_k-\alpha|<\frac{\epsilon}{2}
    \]
    が成り立つ。

    ここで，有限個の和
    \[
        C=
        \sum_{k=1}^{N_1-1}k|a_k-\alpha|
    \]
    を考える。
    これは有限個の実数の和であるから，ある実数として定まる。
    また，各項は $0$ 以上なので
    \[
        C\geq 0
    \]
    である。
    $n\geq N_1$ のとき，
    \[
        |R_n|
        =
        \left|
            \frac{1}{n^2}\sum_{k=1}^{n}k(a_k-\alpha)
        \right|
    \]
    である。

    三角不等式より，
    \[
        |R_n|
        \leq
        \frac{1}{n^2}\sum_{k=1}^{n}k|a_k-\alpha|
    \]
    である。

    この和を
    \[
        \sum_{k=1}^{n}k|a_k-\alpha|
        =
        \sum_{k=1}^{N_1-1}k|a_k-\alpha|
        +
        \sum_{k=N_1}^{n}k|a_k-\alpha|
    \]
    と分ける。

    したがって，
    \[
        |R_n|
        \leq
        \frac{1}{n^2}
        \sum_{k=1}^{N_1-1}k|a_k-\alpha|
        +
        \frac{1}{n^2}
        \sum_{k=N_1}^{n}k|a_k-\alpha|
    \]
    である。

    先ほどの $C$ を用いれば，
    \[
        \frac{1}{n^2}
        \sum_{k=1}^{N_1-1}k|a_k-\alpha|
        =
        \frac{C}{n^2}
    \]
    である。
    また，$k\geq N_1$ ならば
    \[
        |a_k-\alpha|<\frac{\epsilon}{2}
    \]
    であるから，
    \[
        \sum_{k=N_1}^{n}k|a_k-\alpha|
        <
        \sum_{k=N_1}^{n}k\frac{\epsilon}{2}
        \leq
        \frac{\epsilon}{2}\sum_{k=1}^{n}k
    \]
    である。

    よって
    \[
        \frac{1}{n^2}
        \sum_{k=N_1}^{n}k|a_k-\alpha|
        \leq
        \frac{\epsilon}{2}
        \cdot
        \frac{1}{n^2}\sum_{k=1}^{n}k
    \]
    となる。

    さらに
    \[
        \frac{1}{n^2}\sum_{k=1}^{n}k
        =
        \frac{n(n+1)}{2n^2}
        =
        \frac{n+1}{2n}
    \]
    である。
    $n\geq 1$ ならば
    \[
        \frac{n+1}{2n}\leq 1
    \]
    である。
    したがって，
    \[
        \frac{1}{n^2}
        \sum_{k=N_1}^{n}k|a_k-\alpha|
        \leq
        \frac{\epsilon}{2}
    \]
    である。

    以上より，$n\geq N_1$ のとき
    \[
        |R_n|
        \leq
        \frac{C}{n^2}
        +
        \frac{\epsilon}{2}
    \]
    を得る。
    次に，$C/n^2$ を小さくする。
    もし $C=0$ ならば
    \[
        \frac{C}{n^2}=0<\frac{\epsilon}{2}
    \]
    である。

    もし $C>0$ ならば，アルキメデスの性質より，
    \[
        N_2>\sqrt{\frac{2C}{\epsilon}}
    \]
    を満たす自然数 $N_2$ を取ることができる。
    すると，$n\geq N_2$ ならば
    \[
        n^2\geq N_2^2>\frac{2C}{\epsilon}
    \]
    であるから，
    \[
        \frac{C}{n^2}<\frac{\epsilon}{2}
    \]
    となる。

    したがって，いずれの場合も，ある自然数 $N_2$ が存在して，
    $n\geq N_2$ ならば
    \[
        \frac{C}{n^2}<\frac{\epsilon}{2}
    \]
    が成り立つ。
    そこで
    \[
        N=\max\{N_1,N_2\}
    \]
    とおく。

    このとき，$n\geq N$ ならば
    \[
        n\geq N_1,
        \qquad
        n\geq N_2
    \]
    である。

    したがって，
    \[
        |R_n|
        \leq
        \frac{C}{n^2}
        +
        \frac{\epsilon}{2}
        <
        \frac{\epsilon}{2}
        +
        \frac{\epsilon}{2}
        =
        \epsilon
    \]
    となる。

    よって
    \[
        R_n\to 0
    \]
    である。
    以上より，
    \[
        x_n
        =
        \alpha\frac{1}{n^2}\sum_{k=1}^{n}k
        +
        R_n
    \]
    であり，
    \[
        \alpha\frac{1}{n^2}\sum_{k=1}^{n}k
        \to
        \frac{\alpha}{2},
        \qquad
        R_n\to 0
    \]
    である。

    したがって
    \[
        x_n\to \frac{\alpha}{2}
    \]
    である。

    ゆえに
    \[
        \lim_{n\to\infty}
        \frac{a_1+2a_2+\cdots+(n-1)a_{n-1}+na_n}{n^2}
        =
        \frac{\alpha}{2}
    \]
    である。
\end{proof}